% ======================================================================
% SECTION 3: HR 4501 - Patient Self-Selection Act of 2026
% Adaption of 21 CFR Part 50, FDA Decentralized Trial Final Guidance
% 2024, and 42 USC 300gg-8.
%
% Layout follows the prior site documentation template at
% new-trial/site/all-documents/all\_documents.tex (Findings,
% Definitions, Authorization or Rights, Implementation, Reporting and
% Enforcement). The section title in main.tex already names the bill
% as an Adaption.
% ======================================================================

[PROCESSING INSTRUCTION (HR 4501 - Patient Self-Selection Act of 2026,
6 to 9 paragraphs structured into the five legislative-act
subsections): Draft the proposed bill text in legislative-act voice,
following the layout used by SB 1042 in the prior site template at
new-trial/site/all-documents/all\_documents.tex (which was, in turn,
populated from new-trial/site/01-legislation-authorization/physical\_ai\_trial\_authorization.tex
and national-platform/new\_trial\_psl/01\_preamble\_and\_sb1042.tex).

The bill must explicitly state in its first paragraph that it is an
"Adaption of 21 CFR Part 50, FDA Decentralized Trial Final Guidance
2024, and 42 USC 300gg-8" so that readers can trace the chain of prior
authority. Cite \cite{hhs_45cfr46_protection_human_subjects},
\cite{fda_dct_guidance_2024}, \cite{fda_oce_advancing_oncology_dct_2024},
and \cite{usc_300gg8_clinical_trials}.

Subsection 3.1 - Findings and Declarations: Read
patients/paper/Deep-Research-3/part1\_legal\_baseline.md and quote the
seven percent participation-rate finding. Read
patients/paper/Deep-Research-3/part2\_ai\_robotics\_legislation.md and
quote the TrialGPT, PRISM, and TrialMatchAI performance figures.
Compose 4 to 6 enumerated findings in the (a)/(b)/(c) style used by
SB 1042 in the prior template. Each finding must cite at least one
prior-law reference and at least one AI-evidence reference. Cite
\cite{jin2024trialgpt}, \cite{gupta2024prism}, and
\cite{abdallah2026trialmatchai}.

Subsection 3.2 - Definitions: Define the following operational terms
used throughout the bill:
- "Patient self-selection" means the patient's right to discover,
  evaluate, and book their own oncology clinical trial slot without
  intermediary provider gatekeeping. Reference the existing 24/7 trial
  format documented in
  new-trial/national-24-7-trial/paper/full-paper/main.tex and the
  hour-NN/patient\_records.md files at
  new-trial/national-24-7-trial/hour-00/ through hour-55/.
- "Machine-readable eligibility report" means a structured FHIR-based
  output containing the patient's eligibility-relevant data fields
  for at least 80 percent of currently recruiting U.S. oncology
  trials. Reference the ClinicalTrials.gov infrastructure under FDAAA
  Section 801. Cite \cite{rank5_fdaaa_2007}.
- "Qualified Physical AI oncology trial site" means a site authorized
  under either this Act or under the prior California SB 1042 / Title
  22 framework documented at
  new-trial/site/05-state-regulations/ca\_state\_regulations.tex and
  national-platform/new\_trial\_psl/05\_title22\_regulations.tex (read
  this file in full when populating).
- "24/7 booking" means the patient's right to book any qualified site
  at any hour of the day, any day of the week, drawing on the
  168-hour sponsor extension at
  sponsor/final\_paper/168\_hours/ for the operational baseline. Cite
  \cite{kawchak_2026_19396256} and \cite{kawchak_2026_19994945}.

Subsection 3.3 - Patient Self-Selection Rights (the operative core):
Compose 5 to 7 enumerated rights in the (a)/(b)/(c) style. Each right
must be specific enough to be enforceable. Examples:
- Right to receive a machine-readable eligibility report from any
  CMS-billable EHR within 7 calendar days of pathology or genomic
  result release.
- Right to discover all currently recruiting U.S. oncology trials via
  an FDA-published API that integrates the data fields from
  ClinicalTrials.gov, the FDA RTCT pilot program, and FDA Oncology
  Center of Excellence DCT directory at
  new-trial/national-24-7-trial/FDA-April-2026/FDA\_RealTime\_Clinical\_Trials.md
  and \cite{fda_oce_advancing_oncology_dct_2024}.
- Right to book a 24/7 trial slot at any qualified site without prior
  written referral. Reference the 24/7 site infrastructure at
  new-trial/national-24-7-trial/hour-NN/diagram\_facility.txt files
  (e.g., hour-00, hour-12, hour-23, hour-47) for the operational
  pattern.
- Right to payment parity for decentralized and at-home trial
  components under 42 USC 300gg-8 and CMS NCD 310.1; closes the gap
  identified in the Introduction. Cite
  \cite{usc_300gg8_clinical_trials} and
  \cite{cms_ncd_3101_routine_costs}.
- Right to cross-state portability of eligibility data and trial
  bookings, drawing on the ASTP/ONC Hospital APIs Brief 2024 to
  identify the technical baseline. Cite
  \cite{astp_onc_hospital_apis_2024}.

Subsection 3.4 - Implementation: Implementation timeline (12 months
for federal API, 24 months for full state-level adoption); FDA, ONC,
HHS, and CMS roles; preemption rule that explicitly preserves stronger
state-level patient protections (e.g., CA SB 37 for routine-cost
coverage). Cite \cite{rank8_california_sb37_2001} and
\cite{rank4_fdora_fda_modernization_2022} for the FDORA precedent
that authorizes mandatory FDA decentralized-trial guidance.

Subsection 3.5 - Reporting and Enforcement: Annual public report on
patient-self-selection rates by state, payer, and cancer type;
information-blocking penalties harmonized with
\cite{federal_register_info_blocking_disincentives_2024}; an explicit
"no automated-only denial" clause that mirrors the first guardrail
from patients/paper/Deep-Research-3/part3\_metrics\_guardrails.md.

Close the bill (1 paragraph after Subsection 3.5) with two patient-
control framings:
(i) FOR INDIVIDUAL PATIENTS: An NSCLC patient at home in rural Texas
    can, on their own initiative, receive an FHIR eligibility report
    for the qualified Physical AI Oncology Trial Site documented in
    national-platform/new\_trial\_psl/01\_preamble\_and\_sb1042.tex within
    7 days of their genomic result, then self-book a 24/7 slot at a
    USL-7-or-above surgical robot system following the 10-stage
    journey for PAT-2026-0042 in
    patient-journey/stage\_01\_prescreening.py through
    stage\_10\_closeout.py.
(ii) FOR BROAD ADOPTION: All U.S. oncology trials must publish their
     eligibility criteria in machine-readable form within the
     12-month federal-API window, with ALL CMS-billable EHRs
     maintaining a self-selection portal endpoint accessible by every
     U.S. cancer patient by month 24. The United States must remain
     Number 1 in the world for cancer patient benefit during the
     medical AI and robotics revolution, and HR 4501 is the
     foundational bill that opens that pathway. Cite
     \cite{kawchak_2026_19994945} and \cite{kawchak_2026_19119939}.]
